% *- coding: UTF-8 -*-
% !TEX builder = Texlive
% !TEX prggram = xelatex

%文档设置及格式
\documentclass[UTF8,a4paper,12pt,titlepage]{ctexart}
\include{document_setting}

% 设置文件属性内容
\hypersetup{pdfauthor={吴毓明},
            pdftitle={PAM制备投加装置},
            pdfsubject={使用说明},
            pdfkeywords={任何问题，联系61508712@qq.com}}

% 定义页眉页脚
\include{foot_head}

% 表格格式定义
\include{table_format}

% 图片
\include{picture_configuration}

% 特殊符号(例如千分号等)
\usepackage{textcomp}

% 流程图

% 流程图
% 导入tikz绘图包
\usepackage{tikz}
% 导入基本图形
\usetikzlibrary{shapes,arrows}
% 定义流程图基本元素
    %开始结束
    \tikzstyle{startstop} = [
        rectangle,
        rounded corners,
        minimum width=3cm,
        minimum height=1cm,
        text centered, 
        draw=black,
        fill=red!30 %填充颜色
        ]
    %输入输出
    \tikzstyle{io} = [
        trapezium, 
        trapezium left angle = 70,
        trapezium right angle=110,
        minimum width=3cm,
        minimum height=1cm,
        text centered,
        draw=black,
        fill=blue!30 %填充颜色
        ]
    %操作
    \tikzstyle{process} = [
        rectangle,
        minimum width=2.5cm,
        minimum height=1cm,
        text centered,
        text width =2.5cm,
        draw=black,
        fill=orange!30 %填充颜色
        ]
    %判断
    \tikzstyle{decision} = [
        diamond,
        minimum width=3cm,
        minimum height=1cm,
        text centered,
        draw=black,
        fill=green!30 %填充颜色
        ]
    %箭头
    \tikzstyle{arrow} = [thick,->,>=stealth]
%说明文字
\tikzstyle{text_box} = [
   rectangle,
   rounded corners,
   minimum width=1.5cm,
   minimum height=0.5cm,
   text centered, 
   ]
%箭头
\tikzstyle{arrow} = [thick,->,>=stealth]



% 开始正文
\begin{document}

% 制作封面
\title{
   \zihao{1} 高锰酸钾上料设备\\[2.5mm]
   \zihao{1} 试运行技术方案\\
   {\noindent}\rule{16cm}{1pt}\\[3mm]
   \zihao{2} 泊多科技（广州）有限公司\\
   % \zihao{4} 吴毓明\\
   %\zihao{4} Wym\\
   %\zihao{4} 2022年12月27日
   \begin{figure}[h]
      \centering
      \includegraphics[height=10cm]{p1.PNG}
   \end{figure}
}

\maketitle % 封面



\tableofcontents % 目录

\newpage % 换页

%\part{系统组成}
%ctex的\chapter{sdsg}是不能用的

\section{系统组成及功能说明}
   \subsection{系统主要设备}
   高锰酸钾上料设备 的主要组成部分，如下图\ref{fig:p1}所示。

   % 图片模板
      \begin{figure}[h]
         \centering

         \begin{tikzpicture}
            %定义图像
            \node at (0,0) {\includegraphics[height=10cm]{f2.PNG}};

            %说明
            \node at (7,3) (1) [text_box] {1号储液罐};
            \draw [arrow] (1) -- (6,-2);
            \node at (5,4) (2) [text_box] {2号储液罐};
            \draw [arrow] (2) -- (3.5,-0.5);
            \node at (3,5) (3) [text_box] {3号储液罐};
            \draw [arrow] (3) -- (1,1);
            \node at (-2,5) (4) [text_box] {吸料控制柜};
            \draw [arrow] (4) -- (-2.3,1.4);
            \node at (-1,-2) (5) [text_box] {投加机};
            \draw [arrow] (5) -- (-3.2,1.2);
            \node at (-2,-3) (6) [text_box] {送料风机};
            \draw [arrow] (6) -- (-4.7,-0.6);
            \node at (-4,-3) (7) [text_box] {空压机};
            \draw [arrow] (7) -- (-5,-1.3);
            \node at (-6,-4) (8) [text_box] {主控制柜};
            \draw [arrow] (8) -- (-6.5,-1);
            \node at (-4,4) (9) [text_box] {吸料风机};
            \draw [arrow] (9) -- (-5,-0.4);
            \node at (-6,3) (10) [text_box] {冷干机};
            \draw [arrow] (10) -- (-5.7,-0.7);
            \node at (-7,2) (14) [text_box] {储气罐};
            \draw [arrow] (14) -- (-6.6,0);

			% 辅助线
             %\def \xLimit {8};
             %\def \yLimit {8};
             %\input{auxiliary_line}

         \end{tikzpicture}
         \caption{高锰酸钾上料设备系统组成}\label{fig:p1}
      \end{figure}


   系统主要设备组成及说明如下：
   \begin{itemize}
       \item 吸料控制柜：吸料操作控制，触摸屏显示吸料过程运行状态。
       \item 投加机：高锰酸钾粉体物料 的主要储存、投加、溶配主要设备，详细说明见\ref{main_machine}。
       \item 送料风机：提供动力，将 高锰酸钾粉体物料 从 投加机 输送到 1号、2号、3号储液罐。
       \item 吸料风机：提供动力，将 高锰酸钾粉体物料 吸入到 储存破拱料斗。
       \item 空压机：给整套系统提供压缩空气。
       \item 冷干机：通过冷却的方法，干燥 压缩空气，为系统提供干燥的 压缩空气 源。
           通过向 储存破拱料斗 加入 干燥的 压缩空气，保证 储存破拱料斗内湿度 低于设定值，防止物料结拱结块。
       \item 储气罐：储存 压缩空气。
       \item 主控制柜：系统自动控制主要设备，安装PLC、触摸屏等。
   \end{itemize}

   \subsection{投加机组成说明}\label{main_machine}
   投加机 的主要组成部分，如下图\ref{fig:p2}所示。

   
      \begin{figure}[h]
         \centering

         \begin{tikzpicture}
            %图片
            \node at (0,0) {\includegraphics[height=10cm]{f1.PNG}};

            %说明
            \node at (2,5) (1) [text_box] {吸料料斗};
            \draw [arrow] (1) -- (0,4);
            \node at (-4,4) (2) [text_box] {温湿度探头};
            \draw [arrow] (2) -- (-1.2,2.7);
            \node at (-5,1) (3) [text_box] {储存破拱料斗};
            \draw [arrow] (3) -- (-0.4,0.2);
            \node at (6,0) (4) [text_box] {空气炮};
            \draw [arrow] (4) -- (1.6,0.1);
            \node at (-4,-1) (5) [text_box] {气碟};
            \draw [arrow] (5) -- (0,-1.2);
            \node at (5,-1) (6) [text_box] {出仓阀};
            \draw [arrow] (6) -- (0,-2);
            \node at (5,-2) (7) [text_box] {螺旋进料器};
            \draw [arrow] (7) -- (0.7,-2.4);
            \node at (-4,-2) (8) [text_box] {下料阀};
            \draw [arrow] (8) -- (-1,-3);
            \node at (-5,-4) (9) [text_box] {送料射流器};
            \draw [arrow] (9) -- (-1,-3.5);

			% 辅助线
            %\def \xLimit {8};
            %\def \yLimit {8};
            %\input{auxiliary_line}

         \end{tikzpicture}
         \caption{投加机组成}\label{fig:p2}
      \end{figure}

   投加机主要设备组成及说明如下：
   \begin{itemize}
       \item 吸料料斗：暂存吸料机吸入的 高锰酸钾粉体物料。
       \item 温湿度探头：检测 储存破拱料斗内湿度。若湿度高于设定值，则向 储存破拱料斗内 通入 干燥冷空气，防止物料结拱结块。
       \item 储存破拱料斗：储存 高锰酸钾粉体物料。
       \item 空气炮：当 储存破拱料斗内 湿度过高 或 投加过程中，向 储存破拱料斗内 通入 高压 干燥冷空气，防止物料结拱结块。
       \item 气碟：系统投加溶配工作时，辅助 储存破拱料斗内 高锰酸钾粉体物料 流动。
       \item 出仓阀：控制 高锰酸钾粉体物料 进入 螺旋进料器。
       \item 螺旋进料器：定量输送 高锰酸钾粉料， 确保投加量准确。
       \item 下料阀：控制 高锰酸钾粉体物料 进入 送料射流器。
   \end{itemize}

\section{系统启动}\label{sec:sg1}
   \subsection{检查空压机油位}
   系统启动前，先检查 空压机 气液分离器油位。
   气液分离器 如下图\ref{fig:p3}所示。
   % 图片模板
      \begin{figure}[h]
         \centering

         \begin{tikzpicture}
            %定义图像
            \node at (0,0) {\includegraphics[height=10cm]{f3.PNG}};

            %说明
            \def \itemNumber {1};
            \def \centertPoint {0,4};
            \node at (-5,4) (\itemNumber) [text_box] {加油口};
            \draw [arrow] (\itemNumber) -- (\centertPoint);
            \draw (\centertPoint) circle (0.7)[red];
            \def \itemNumber {2};
            \def \centertPoint {0.2, -2.5};
            \node at (5,-2) (\itemNumber) [text_box] {油位视窗};
            \draw [arrow] (\itemNumber) -- (\centertPoint);
            %\draw (\centertPoint) circle (0.5)[red];
            \def \itemNumber {3};
            \def \centertPoint {0, -4.5};
            \node at (-5,-4) (\itemNumber) [text_box] {排水口};
            \draw [arrow] (\itemNumber) -- (\centertPoint);
            \draw (\centertPoint) circle (0.7)[red];

			% 辅助线
             \def \xLimit {8};
             \def \yLimit {8};
             %\input{auxiliary_line}

         \end{tikzpicture}
         \caption{空压机 气液分离器}\label{fig:p3}
      \end{figure}


   观察 油位视窗，如果 油位过低，则从 加油口 加入 46号螺杆空压机润滑油；
   如果 在 油位视窗 中发现 润滑油含水量过多，则打开下部 排水口，排出 气液分离器内水份。

   \subsection{检查系统管道系统}
   检查 高锰酸钾上料设备系统 的物料管路、吸风管路、压缩空气管路，确保管路通畅。

   \subsection{系统上电}\label{sec:sg2}
   依次接通上级电柜中 高锰酸钾上料设备 电源开关，以及 主控制柜 中 总电源开关。
   
        \begin{figure}[h]
            \centering
            \begin{tikzpicture}
                %图片
                \node at (-3,0) {\includegraphics[height=10cm]{f4.PNG}};
                \node at (5,0) {\includegraphics[height=10cm]{f5.PNG}};

                %说明
                \def \itemNumber {1};
                \def \centertPoint {4.6, -1.5};
                \node at (5,-5.5) (\itemNumber) [text_box] {上级电柜-本系统电源开关};
                \draw [arrow] (\itemNumber) -- (\centertPoint);
                \draw (\centertPoint) circle (0.7)[red];
                \def \itemNumber {2};
                \def \centertPoint {-5.5,2.3};
                \node at (-4,-5.5) (\itemNumber) [text_box] {主控制柜-总电源开关};
                \draw [arrow] (\itemNumber) -- (\centertPoint);
                \draw (\centertPoint) circle (0.8)[red];

                % 辅助线
                \def \xLimit {8};
                \def \yLimit {8};
                %\input{auxiliary_line}

            \end{tikzpicture}
            \caption{电源开关位置}\label{fig:p4}
        \end{figure}


   \par系统上电后，空压机 会自动启动，直到 压缩空气压力达到0.5-0.7MPa时，才能进行吸料操作或投加操作。
   压缩空气压力未达设置，系统会“气压低”报警，无法 进行吸料操作或投加操作。
   \par两个开关位置 如下图\ref{fig:p4}所示。

\section{吸料系统操作规程}
   \subsection{吸料工作流程}
   吸料工作流程 如下：
   \begin{enumerate}
       \item 启动吸料。
       \item 系统自动通过压缩空气反冲清理吸料系统滤芯。
       \item 启动 吸料风机，将 高锰酸钾粉料 吸入到 吸料料斗。
       \item 到达设定 吸料时间，停止 吸料风机，高锰酸钾粉料 从 吸料料斗 落入 储存破拱料斗。
       \item 系统自动通过压缩空气反冲清理吸料系统滤芯。
       \item 启动 吸料风机，将 高锰酸钾粉料 吸入到 吸料料斗。
       \item 上述两步循环运行，直到吸完所有物料，停止吸料。
   \end{enumerate}

   \newpage

   \subsection{吸料控制柜操作界面}
   
         \begin{figure}[h]
            \centering
            \begin{tikzpicture}
               %图片
               \node at (0,0) {\includegraphics[height=7cm]{p11.PNG}};

               %说明
               %\def \itemNumber {1};
               %\def \centertPoint {-1.9, 2.7};
               %\node at (-3,3) (\itemNumber) [text_box] {电柜总开关};
               %\draw (\centertPoint) rectangle +(0.8,0.9)[red];
               %\def \itemNumber {2};
               %\def \centertPoint {-1.9, 1.5};
               %\node at (-3,2) (\itemNumber) [text_box] {其余开关};
               %\draw (\centertPoint) rectangle +(3.2,0.7)[red];

            % 辅助线
               \def \xLimit {8};
               \def \yLimit {7};
            %    \input{auxiliary_line}

            \end{tikzpicture}
            \caption{吸料控制柜操作界面}\label{fig:p5}
         \end{figure}



   \subsection{吸料参数设置}
   \input{picture6}
   \par吸料参数 有 吸料时间、落料延时、滤芯清理次数3项参数设置：
   \begin{itemize}
       \item 吸料时间：吸料风机每次运行的时间，建议设置为1次或2次吸完1罐物料较有合适，具体时间建议由实验测定。
       \item 落料延时：高锰酸钾粉体物料 从 吸料料斗 落入 储存破拱料斗用时，此时间内不启动 吸料风机 也不 清理吸料系统滤芯。
       \item 滤芯清理次数：吸料系统启动前或吸料风机运行间隙，吸料系统自动清理滤芯的次数。
   \end{itemize}

   \newpage

   \subsection{吸料操作}
   

\begin{figure}[h]
 \centering
    \begin{tikzpicture}
       %图片
       \node at (0,0) {\includegraphics[height=6cm]{f6.PNG}};

        %说明
        \def \itemNumber {1};
        \def \centertPoint {-0.3,-0.6};
        \node at (-5,0) (\itemNumber) [text_box] {启动按钮};
        \draw [arrow] (\itemNumber) -- (\centertPoint);
        \draw (\centertPoint) circle (0.3)[red];
        \def \itemNumber {2};
        \def \centertPoint {0.3,-0.6};
        \node at (5,0) (\itemNumber) [text_box] {停止按钮};
        \draw [arrow] (\itemNumber) -- (\centertPoint);
        \draw (\centertPoint) circle (0.3)[red];

        % 辅助线
       \def \xLimit {8};
       \def \yLimit {5};
       %\input{auxiliary_line}

    \end{tikzpicture}
 \caption{吸料控制柜}\label{fig:p7}
\end{figure}


   \par按“启动”按钮开始吸料，按“停止”按钮停止吸料。 

\section{投加系统操作规程}
   \subsection{主操作界面}
   \input{picture8}
   \newpage

   \subsection{系统参数设置}
   系统参数设置 主要有 空压机、冷干机运行及报警参数设置。
   
          \begin{figure}[h]
             \centering
               \begin{tikzpicture}
                  %图片
                  \node at (0,0) {\includegraphics[height=7cm]{p7.PNG}};

                  %说明

            % 辅助线
               \def \xLimit {8};
               \def \yLimit {4};
                %\input{auxiliary_line}

               \end{tikzpicture}
            \caption{空压机参数设置}\label{fig:p9}
         \end{figure}


   \par空压机参数设置有 压缩空气压力过低设置值1项参数。
   
          \begin{figure}[h]
             \centering
               \begin{tikzpicture}
                  %图片
                  \node at (0,0) {\includegraphics[height=7cm]{p8.PNG}};

                  %说明
                    %\draw (1.7,1.2) rectangle +(1,1)[red];

            % 辅助线
               \def \xLimit {8};
               \def \yLimit {4};
                %\input{auxiliary_line}

               \end{tikzpicture}
            \caption{冷干机参数设置}\label{fig:p10}
         \end{figure}


   \par冷干机参数设置有 开机湿度、停机湿度、高湿度报警3项参数设置。
   \newpage

   \subsection{投加参数设置}\label{sec:sg3}
   % 图片模板
      \begin{figure}[h]
         \centering

         \begin{tikzpicture}
            %定义图像
            \node at (0,0) {\includegraphics[height=7cm]{p9.PNG}};

            %说明
            %\def \itemNumber {1};
            %\def \centertPoint {-3.9,1.1};
            %\node at (-7,3) (\itemNumber) [text_box] {检修旁通阀};
            %\draw [arrow] (\itemNumber) -- (\centertPoint);
            %\draw (\centertPoint) circle (0.5)[red];
            %\def \itemNumber {2};
            %\def \centertPoint {-0.3, -4.2};
            %\node at (-2,-6) (\itemNumber) [text_box] {放空阀};
            %\draw [arrow] (\itemNumber) -- (\centertPoint);
            %\draw (\centertPoint) circle (0.5)[red];

			% 辅助线
             \def \xLimit {8};
             \def \yLimit {8};
             %\input{auxiliary_line}

         \end{tikzpicture}
         \caption{投加机操作界面}\label{fig:p11}
      \end{figure}


   投加参数可设置如下：
   \begin{itemize}
       \item 重量过高报警：储存破拱料斗内 存料重量过高时报警的报警值。
       \item 重量过低报警：储存破拱料斗内 存料重量过低时报警的报警值。
       \item 投加点选择：可设置1号储液罐、2号储液罐、3号储液罐。
       \item 投加重量设置：设置本次投加的重量，单位kg。
   \end{itemize}

   \subsection{自动运行操作}
   按\ref{sec:sg3}说明在图\ref{fig:p11}设置好投加参数之后，点击“开始投加”，系统会自动按照图\ref{fig:p11}投加流程投加。
   系统投加 高锰酸钾粉料 至设定重量后，会自动按投加流程停止。
   投加完成后，送料风机 会继续运行1分钟，吹干净管道系统内残余物料。
   \par 投加过程中，需要中途停止投加，可以按“投加中止”按钮，中止投加。

   \subsection{投加工作流程}\label{sec:sg4}
   投加 高锰酸钾粉料 的 工作流程如下：
   \begin{enumerate}
       \item 启动送料风机。
       \item 开下料阀：系统检测各压力探头压力值，达到设定范围，打开下料阀。
       \item 启动螺旋进料器。
       \item 开出仓阀。
       \item 向储液罐输送 高锰酸钾粉料。输送过程中，系统检测各传感器参数，有异常变化会自己运行清理程序防止堵料。
       \item 关出仓阀：接近目标投加重量时，关出仓阀，但暂不停止 螺旋进料器 和关闭下料阀，将 螺旋进料器 内物料排净。
       \item 停止螺旋进料器。
       \item 关下料阀：关下料阀 后，暂不停止 送料风机，将管道内残余物料吹净。
       \item 停止风机。
   \end{enumerate}

   \subsection{手动运行操作}
   如果需要 手动运行操作，将 主控制柜 上 “手动/自动”切换开关调到“手动”。
   然后按上述\ref{sec:sg4}流程，在 主控制柜 操作面板上，手动操作各设备运行、停止。
   \input{picture12}

\section{停机注意事项}
   \begin{itemize}
       \item 当系统所有设备停止运行后，才切断\ref{sec:sg2}中两个电源开关。
       \item 当 储存破拱料斗内 仍然存放 高锰酸钾粉体物料 时，不建议切换系统电源。
           因为环境湿度较高时，高锰酸钾粉体物料 极容易 结拱结块、堵塞设备及管道，
           需要根据 温湿度探头 反馈向 储存破拱料斗内 通入 干燥冷空气，保持 储存破拱料斗内 干燥。
   \end{itemize}


\end{document}
